\documentclass[12pt,a4paper]{article}

\usepackage[a4paper, margin=1in]{geometry}
\usepackage{mathptmx}                 % Times New Roman
\usepackage{amsmath, amssymb, amsthm}
\usepackage{mathtools}
\usepackage{enumitem}
\usepackage{xcolor}
\usepackage{tcolorbox}
\usepackage{titlesec}
\usepackage{hyperref}
\hypersetup{colorlinks=true, linkcolor=blue!50!black, urlcolor=blue!50!black}

\titleformat{\section}{\normalfont\Large\bfseries}{\thesection.}{0.5em}{}
\titleformat{\subsection}{\normalfont\large\bfseries}{\thesubsection)}{0.5em}{}

\newtcolorbox{insight}{
  colback=blue!4, colframe=blue!40!black, boxrule=0.5pt,
  left=8pt, right=8pt, top=4pt, bottom=4pt,
  before skip=8pt, after skip=8pt
}
\newtcolorbox{warn}{
  colback=red!4, colframe=red!50!black, boxrule=0.5pt,
  left=8pt, right=8pt, top=4pt, bottom=4pt,
  before skip=8pt, after skip=8pt
}
\newtcolorbox{sanity}{
  colback=green!4, colframe=green!40!black, boxrule=0.5pt,
  left=8pt, right=8pt, top=4pt, bottom=4pt,
  before skip=8pt, after skip=8pt
}

\newcommand{\M}{\mathbb{M}}
\newcommand{\vnu}{\vec{\nu}}
\newcommand{\vP}{\vec{P}}
\renewcommand{\d}{\,\mathrm{d}}

\title{\bfseries SIMCC 2026 --- Learning Guide for Questions 1 \& 2\\ \large Markov Matrices, the Power Method, and SVD}
\author{Walkthrough of parts 1a -- 1h and 2d, 2e}
\date{}

\begin{document}
\maketitle

\section*{How to read this guide}

This is a study companion, not a model answer. For each sub-task you get: (i)~what the question is really asking, (ii)~the mathematical derivation step by step, (iii)~the conceptual insight that earns interpretation marks, and (iv)~common pitfalls. The actual report should be \emph{much shorter}: aim for a few clean lines per part, not pages. The full report is capped at 10 pages and Question 3 should get the lion's share of space.

\section*{Setup recap}

Two competing species with fractional populations $p_1(t), p_2(t) \in [0,1]$, evolved by
\[
\begin{pmatrix} p_1(t+1) \\ p_2(t+1) \end{pmatrix}
= \underbrace{\begin{pmatrix} 1-\epsilon & 0 \\ \epsilon & 1 \end{pmatrix}}_{\M}
\begin{pmatrix} p_1(t) \\ p_2(t) \end{pmatrix}, \qquad 0 < \epsilon < 1.
\tag{1}
\]
In plain English: at each generation, a fraction $\epsilon$ of type-1 organisms convert into type-2 organisms, and conversions only go one way (no type-2 ever turns back into type-1).

\begin{insight}
\textbf{Two structural facts about }$\M$ \textbf{that drive the entire question:}
\begin{enumerate}[label=(\arabic*),leftmargin=*]
    \item Each \emph{column} of $\M$ sums to $1$: column 1 is $(1-\epsilon) + \epsilon = 1$, column 2 is $0 + 1 = 1$. This is the defining property of a (column-stochastic) Markov matrix, and it is what makes total population conserved (part 1a).
    \item $\M$ is \emph{lower triangular}. The eigenvalues of any triangular matrix are exactly its diagonal entries, so before computing anything we already know they will be $\{1,\, 1-\epsilon\}$ (part 1b). The derivation is still expected, but the answer should not surprise you.
\end{enumerate}
\end{insight}

\section*{1a) Conservation [1 point]}

\paragraph{Goal.} Show that $p_1(t) + p_2(t)$ does not change with $t$.

\paragraph{Proof (direct algebra).} From the recursion,
\[
p_1(t+1) + p_2(t+1) \;=\; (1-\epsilon)\,p_1(t) \;+\; \bigl[\epsilon\, p_1(t) + p_2(t)\bigr] \;=\; p_1(t) + p_2(t).
\]
So the sum is the same one generation later. By induction on $t$, it is the same at every generation:
\[
p_1(t) + p_2(t) \;=\; p_1(0) + p_2(0) \quad \text{for all } t \ge 0.
\]

\paragraph{Why it works (cleaner argument).} For any column-stochastic matrix $\M$, the row vector $\vec{1}^{\top} = (1,1)$ satisfies $\vec{1}^{\top}\M = \vec{1}^{\top}$, because $\vec{1}^{\top}\M$ just sums the columns. Hence
\[
p_1(t+1) + p_2(t+1) \;=\; \vec{1}^{\top}\vP(t+1) \;=\; \vec{1}^{\top}\M\,\vP(t) \;=\; \vec{1}^{\top}\vP(t) \;=\; p_1(t) + p_2(t).
\]
This is the same statement, but it makes the role of the column-sum-1 property visible. In the report, either argument is fine; the second is slightly more general and shorter.

\section*{1b) Eigenspectra [2 points]}

\paragraph{Goal.} Find the eigenvalues $\lambda_1, \lambda_2$ and eigenvectors $\vnu_1, \vnu_2$ of $\M$.

\paragraph{Step 1 --- characteristic polynomial.}
\[
\M - \lambda I \;=\; \begin{pmatrix} 1-\epsilon-\lambda & 0 \\ \epsilon & 1-\lambda \end{pmatrix},
\quad
\det(\M - \lambda I) \;=\; (1-\epsilon-\lambda)(1-\lambda) - 0\cdot\epsilon \;=\; (1-\epsilon-\lambda)(1-\lambda).
\]
Setting the determinant to zero yields two roots immediately:
\[
\boxed{\lambda_1 = 1, \qquad \lambda_2 = 1-\epsilon.}
\]
Ordering by magnitude: since $0 < \epsilon < 1$, we have $1 > 1-\epsilon > 0$, so $\lambda_1$ is \emph{dominant}. This ordering matters for parts 1d--1e.

\paragraph{Step 2 --- eigenvector for $\lambda_1 = 1$.} Solve $(\M - I)\vnu_1 = \vec{0}$:
\[
\begin{pmatrix} -\epsilon & 0 \\ \epsilon & 0 \end{pmatrix}\begin{pmatrix} a\\b \end{pmatrix} = \begin{pmatrix} 0\\0 \end{pmatrix}
\;\Longrightarrow\; -\epsilon a = 0 \;\Rightarrow\; a = 0,\; b \text{ free}.
\]
Taking $b=1$,
\[
\boxed{\vnu_1 = \begin{pmatrix} 0 \\ 1 \end{pmatrix}.}
\]
\begin{sanity}
\textbf{Sanity check.} $\M\vnu_1 = \begin{pmatrix} (1-\epsilon)\cdot 0 + 0\cdot 1 \\ \epsilon\cdot 0 + 1\cdot 1\end{pmatrix} = \begin{pmatrix} 0\\1\end{pmatrix} = 1\cdot\vnu_1$. \checkmark
\end{sanity}

\paragraph{Step 3 --- eigenvector for $\lambda_2 = 1-\epsilon$.} Solve $(\M - (1-\epsilon)I)\vnu_2 = \vec{0}$:
\[
\begin{pmatrix} 0 & 0 \\ \epsilon & \epsilon \end{pmatrix}\begin{pmatrix} a\\b \end{pmatrix} = \begin{pmatrix} 0\\0 \end{pmatrix}
\;\Longrightarrow\; \epsilon a + \epsilon b = 0 \;\Rightarrow\; b = -a.
\]
Taking $a=1$,
\[
\boxed{\vnu_2 = \begin{pmatrix} 1 \\ -1 \end{pmatrix}.}
\]
\begin{sanity}
\textbf{Sanity check.} $\M\vnu_2 = \begin{pmatrix} (1-\epsilon)\cdot 1 + 0\cdot(-1) \\ \epsilon\cdot 1 + 1\cdot(-1)\end{pmatrix} = \begin{pmatrix} 1-\epsilon\\ \epsilon-1\end{pmatrix} = (1-\epsilon)\begin{pmatrix} 1\\-1\end{pmatrix} = \lambda_2\vnu_2$. \checkmark
\end{sanity}

\paragraph{Interpretation (worth one mark).}
$\vnu_1 = (0,1)^{\top}$ is the all-type-2 state --- a genuine, physically realisable population. It is the system's fixed point, since once everyone is type-2 no further conversion can occur.

$\vnu_2 = (1,-1)^{\top}$ has a \emph{negative} component, so it is \emph{not} a physically realisable population on its own. It is a ``difference mode'': any deviation of the initial state from the steady state $\vnu_1$ is captured by $\vnu_2$ and decays away at rate $(1-\epsilon)^t$. Recognising this distinction is what separates a mechanical answer from an explained one.

\begin{warn}
\textbf{Pitfall.} Do not normalise the eigenvectors to unit length unless the question asks for it; the algebra in 1c stays cleaner with $\vnu_1 = (0,1)^{\top}$ and $\vnu_2 = (1,-1)^{\top}$. Any nonzero scalar multiple of an eigenvector is still an eigenvector.
\end{warn}

\section*{1c) Linear combinations [1 point]}

\paragraph{Setup.} We are told $p_1(0) = \tfrac{1}{2}$. Combined with the conservation law from 1a and the convention that the fractions sum to $1$, we have $p_2(0) = \tfrac{1}{2}$, so
\[
\vP(0) = \begin{pmatrix} 1/2 \\ 1/2 \end{pmatrix}.
\]

\paragraph{Goal.} Find scalars $c_1, c_2$ with $\vP(0) = c_1\vnu_1 + c_2\vnu_2$.

\paragraph{Solve componentwise.}
\[
\begin{pmatrix} 1/2 \\ 1/2 \end{pmatrix} \;=\; c_1\begin{pmatrix} 0\\1\end{pmatrix} + c_2\begin{pmatrix} 1\\-1\end{pmatrix} \;=\; \begin{pmatrix} c_2\\ c_1 - c_2 \end{pmatrix}.
\]
From the first row, $c_2 = \tfrac{1}{2}$. Plugging into the second row, $c_1 - \tfrac{1}{2} = \tfrac{1}{2}$, so $c_1 = 1$.
\[
\boxed{\vP(0) = 1\cdot\vnu_1 + \tfrac{1}{2}\cdot\vnu_2 = \begin{pmatrix} 0\\1\end{pmatrix} + \tfrac{1}{2}\begin{pmatrix} 1\\-1\end{pmatrix}.}
\]

\begin{insight}
\textbf{What's the point of decomposing into eigenvectors?} Because $\M\vnu_i = \lambda_i\vnu_i$, applying $\M$ in this basis is just scalar multiplication of each component by its own $\lambda_i$. So a single iteration of the dynamics becomes
\[
\M\vP(0) \;=\; c_1\lambda_1\vnu_1 + c_2\lambda_2\vnu_2,
\]
and after $t$ iterations the only thing that changes is the exponent on each $\lambda$. This trick collapses the recursion into a closed form, which is what 1d asks for.
\end{insight}

\section*{1d) Eigenvectors make time evolution easy [2 points]}

\paragraph{Goal.} Closed-form expression for $\vP(t)$ in terms of $c_1, c_2, \vnu_1, \vnu_2, \lambda_1, \lambda_2$.

\paragraph{Derivation.} Starting from $\vP(0) = c_1\vnu_1 + c_2\vnu_2$,
\[
\vP(1) = \M\vP(0) = c_1\M\vnu_1 + c_2\M\vnu_2 = c_1\lambda_1\vnu_1 + c_2\lambda_2\vnu_2,
\]
and by induction (apply $\M$ again, scalars factor through),
\[
\boxed{\vP(t) \;=\; c_1\,\lambda_1^{\,t}\,\vnu_1 \;+\; c_2\,\lambda_2^{\,t}\,\vnu_2.}
\]
Substituting our specific values gives
\[
\vP(t) \;=\; 1\cdot 1^{\,t}\begin{pmatrix} 0\\1\end{pmatrix} + \tfrac{1}{2}(1-\epsilon)^{\,t}\begin{pmatrix} 1\\-1\end{pmatrix} \;=\; \begin{pmatrix} \tfrac{1}{2}(1-\epsilon)^{t} \\ 1 - \tfrac{1}{2}(1-\epsilon)^{t}\end{pmatrix}.
\]
\begin{sanity}
\textbf{Sanity check at $t=0$:} $\vP(0) = (1/2,\, 1/2)^{\top}$. \checkmark\quad
Conservation: $p_1(t)+p_2(t) = 1$ for every $t$. \checkmark
\end{sanity}

\section*{1e) Long-time behaviour [2 points]}

\paragraph{Part (i): limit.} Since $0 < 1-\epsilon < 1$, the geometric factor $(1-\epsilon)^t \to 0$ as $t \to \infty$. Hence the $\vnu_2$ component of $\vP(t)$ vanishes and
\[
\vP(t) \;\xrightarrow{\;t\to\infty\;}\; c_1\,\vnu_1 \;=\; \begin{pmatrix} 0\\1\end{pmatrix}.
\]
The population settles into the eigenvector of $\M$ with the larger eigenvalue, $\lambda_1 = 1$.

\paragraph{Part (ii): structural reason.} Two complementary explanations, either is acceptable:
\begin{itemize}[leftmargin=*]
\item \emph{Algebraic.} The recursion in the eigenbasis is $c_i(t) = \lambda_i^{\,t}\,c_i(0)$. Any eigenvalue with $|\lambda|<1$ shrinks each generation, while $\lambda = 1$ is fixed. After many iterations only the $|\lambda|=1$ modes survive --- in this problem there is just one, namely $\vnu_1$.
\item \emph{Physical.} The matrix has zero in the $(1,2)$ position, i.e.~no type-2 organism ever becomes type-1. Every generation drains a fraction $\epsilon$ of the remaining type-1 population into type-2 and never reverses. Type-1 therefore decays geometrically and the system ends up entirely in type-2: that is $\vnu_1$.
\end{itemize}

\begin{insight}
\textbf{The big idea of the whole question.} ``Find the eigenvector with the largest eigenvalue'' is the same as ``find the long-time fate of the system''. This is exactly the principle that Question~2 abstracts into the Power Method and that Question~3 uses on real spectroscopic data: the dominant eigenvector is what survives.
\end{insight}

\section*{1f) Numerical validation of parts (d) and (e) [2 points]}

\paragraph{Goal.} Confirm the closed-form expression for $\vP(t)$ matches a direct simulation, and that $\vP(t) \to (0,1)^{\top}$.

\paragraph{What to compute.}
\begin{enumerate}[label=(\arabic*),leftmargin=*]
    \item Pick a representative $\epsilon$ (e.g.~$\epsilon = 0.1$). Initial state $\vP(0) = (1/2,\, 1/2)^{\top}$.
    \item \textbf{Direct simulation:} iterate $\vP(t+1) = \M\vP(t)$ for $t = 0,\dots,T$ (e.g.~$T=100$).
    \item \textbf{Analytical:} use the closed form $\vP(t) = c_1\vnu_1 + c_2(1-\epsilon)^t\vnu_2$ with $c_1 = 1, c_2 = 1/2$.
    \item Confirm the two agree to machine precision at every $t$ (max absolute error $< 10^{-12}$).
    \item Plot $p_1(t)$ and $p_2(t)$ versus $t$ on a single panel. $p_1$ should decay from $1/2$ to $0$, $p_2$ should grow from $1/2$ to $1$.
\end{enumerate}

\paragraph{Suggested figure for the report.} Two curves on a single plot, log-$y$ axis for $p_1$ would show its geometric decay especially cleanly; a dashed horizontal line at $1$ marks the asymptote of $p_2$.

\begin{warn}
\textbf{Pitfall.} Don't simulate by computing $\M^t$ as a matrix power and then multiplying. For $t\sim 100$ it works, but it is slower and less stable than iterating $\vP \leftarrow \M\vP$. Iterate.
\end{warn}

\section*{1g) Time-varying conversion rate [2 points]}

\paragraph{Setup.} Now $\epsilon(t) = \epsilon_0\,e^{-t/t_0}$ decays with $t$, so the matrix is $\M(t) = \begin{pmatrix} 1-\epsilon(t) & 0\\ \epsilon(t) & 1\end{pmatrix}$ and the recursion is $\vP(t+1) = \M(t)\,\vP(t)$.

\paragraph{Key observation.} The eigenvectors $\vnu_1 = (0,1)^{\top}$ and $\vnu_2 = (1,-1)^{\top}$ \emph{do not depend on $\epsilon$}, so they are the same at every $t$. Only the eigenvalues change: at time step $t$ they are $\lambda_1(t) = 1$ and $\lambda_2(t) = 1 - \epsilon(t)$.

\paragraph{Time evolution.} In the (constant) eigenbasis, write $\vP(t) = c_1(t)\vnu_1 + c_2(t)\vnu_2$. Then
\[
c_1(t+1) = 1\cdot c_1(t) \quad\text{(constant)}, \qquad c_2(t+1) = \bigl(1-\epsilon(t)\bigr)\,c_2(t).
\]
Iterating the second gives a \emph{product}, not a power:
\[
c_2(t) \;=\; c_2(0)\,\prod_{s=0}^{t-1}\bigl(1-\epsilon(s)\bigr).
\]
Hence the closed form for part (g) is
\[
\boxed{\;\vP(t) \;=\; c_1\,\vnu_1 \;+\; c_2\left[\prod_{s=0}^{t-1}\bigl(1-\epsilon_0\,e^{-s/t_0}\bigr)\right]\vnu_2.\;}
\]
The change from part (d) is exactly the replacement of the geometric factor $(1-\epsilon)^t$ by the product $\prod_{s=0}^{t-1}(1-\epsilon(s))$.

\paragraph{The limit $t\to\infty$ (where the twist lives).} For small $\epsilon(s)$, $\ln(1-\epsilon(s)) \approx -\epsilon(s)$, so
\[
\ln\!\prod_{s=0}^{\infty}\bigl(1-\epsilon(s)\bigr) \;\approx\; -\sum_{s=0}^{\infty}\epsilon_0\,e^{-s/t_0} \;=\; -\,\frac{\epsilon_0}{1 - e^{-1/t_0}} \;\approx\; -\epsilon_0\,t_0 \;\text{(for } t_0\gg 1\text{)}.
\]
This sum is \emph{finite}, so the infinite product converges to a finite, \textbf{non-zero} value:
\[
\prod_{s=0}^{\infty}(1-\epsilon(s)) \;\to\; e^{-\epsilon_0 t_0} \;\;(\text{approximately, for } t_0 \gg 1).
\]
\begin{insight}
\textbf{The interesting consequence.} With constant $\epsilon$, the system always reaches the all-type-2 steady state. With $\epsilon(t)$ \emph{decaying fast enough}, the conversion machinery shuts off before all type-1 organisms are converted, so the system gets ``frozen'' at an intermediate state with both species present:
\[
\vP(\infty) \;=\; c_1\,\vnu_1 \;+\; c_2\, e^{-\epsilon_0 t_0}\,\vnu_2 \;\neq\; \vnu_1.
\]
The decay rate competing with the conversion rate is the physics: if $\epsilon_0 t_0$ is large (slow decay), the survivor population is small; if $\epsilon_0 t_0$ is small (fast decay), most type-1 survives.
\end{insight}

\section*{1h) Numerical validation of part (g) [2 points]}

\paragraph{Goal.} Same as 1f, but with the time-varying $\epsilon(t) = \epsilon_0\,e^{-t/t_0}$.

\paragraph{What to compute.}
\begin{enumerate}[label=(\arabic*),leftmargin=*]
    \item Pick representative $\epsilon_0$ and $t_0$ (e.g.~$\epsilon_0 = 0.5$, $t_0 = 10$, so the product limit $e^{-\epsilon_0 t_0} = e^{-5} \approx 0.0067$ is small but nonzero).
    \item \textbf{Direct simulation:} at each step, compute $\epsilon(t)$, build $\M(t)$, and update $\vP(t+1) = \M(t)\vP(t)$.
    \item \textbf{Analytical:} accumulate the product $\Pi(t) = \prod_{s=0}^{t-1}(1-\epsilon(s))$ and compute $\vP(t) = c_1\vnu_1 + c_2\,\Pi(t)\,\vnu_2$.
    \item Verify the two agree to machine precision.
    \item Plot $p_1(t)$ and $p_2(t)$ on the same axes as in 1f, on top of the constant-$\epsilon$ curve from 1f for comparison.
\end{enumerate}

\paragraph{What the plot should show.}
\begin{itemize}[leftmargin=*]
    \item For constant $\epsilon$: $p_1(t) \to 0$, $p_2(t) \to 1$.
    \item For time-decaying $\epsilon(t)$: $p_1(t)$ decays for a while but \emph{plateaus at a positive value} $\tfrac{1}{2}e^{-\epsilon_0 t_0}$ (approximately); $p_2(t)$ correspondingly plateaus below $1$. Show this plateau on a log-$y$ axis.
\end{itemize}

\paragraph{Justification for using this comparison in the report.} The juxtaposition makes the qualitative change visible at a glance: a conversion rate that switches off in finite time leaves a residue of type-1 forever. This is the headline result of 1g and the figure is exactly the right vehicle for the four interpretation marks in 1g+1h.

\section*{Question 2 --- context}

Question 2 (16 points) moves from solving $\det(\M-\lambda I)=0$ by hand to the \emph{Power Method} (an iterative recipe that finds the dominant eigenvector of any symmetric matrix using only matrix--vector products) and the \emph{Singular Value Decomposition} (which extends eigendecomposition to rectangular matrices). Parts 2a--2c prove the Power Method and stress-test it under a vanishing spectral gap; this guide picks up at 2d, where the rectangular-matrix machinery is assembled, and 2e, where the algorithm is committed to pseudocode and verified by hand.

\begin{insight}
\textbf{Why we need }$\mathbb{X}^\top\mathbb{X}$\textbf{ and }$\mathbb{X}\mathbb{X}^\top$. The Power Method's convergence proof in 2a needs an orthonormal eigenbasis, which exists only for symmetric matrices (spectral theorem). Most real data matrices are rectangular ($N\times M$ with $N\ne M$) and have \emph{no eigenvalues at all}: ``$\mathbb{X}\vnu = \lambda\vnu$'' would require $\mathbb{X}\vnu \in \mathbb{R}^N$ to equal $\lambda\vnu \in \mathbb{R}^M$, which is dimensionally impossible. The 2d construction is the standard workaround: from any rectangular $\mathbb{X}$ build two symmetric, positive semi-definite matrices ($\mathbb{X}^\top\mathbb{X}$ in feature space, $\mathbb{X}\mathbb{X}^\top$ in sample space), feed either one to the Power Method, and read singular values and vectors of $\mathbb{X}$ out the other side.
\end{insight}

\section*{2d) Two symmetric matrices: $\mathbb{X}^\top\mathbb{X}$ and $\mathbb{X}\mathbb{X}^\top$ [5 points]}

\paragraph{Goal.} (i) Show $\mathbb{X}^\top\mathbb{X}$ and $\mathbb{X}\mathbb{X}^\top$ are symmetric and positive semi-definite (PSD). (ii) Define singular values $\sigma_i := \sqrt{\lambda_i}$ and prove $\mathbb{X}^\top\mathbb{X}$ and $\mathbb{X}\mathbb{X}^\top$ have the same nonzero eigenvalues, both equal to $\sigma_i^2$.

\paragraph{Part (i) --- symmetry and PSD [2 points].} Two one-line calculations.

\emph{Symmetry.} Using $(AB)^\top = B^\top A^\top$ and $(A^\top)^\top = A$,
\[
(\mathbb{X}^\top\mathbb{X})^\top \;=\; \mathbb{X}^\top (\mathbb{X}^\top)^\top \;=\; \mathbb{X}^\top\mathbb{X}.
\]
Identical algebra gives $(\mathbb{X}\mathbb{X}^\top)^\top = \mathbb{X}\mathbb{X}^\top$.

\emph{PSD.} For any $\vec w \in \mathbb{R}^M$, rewrite the quadratic form as a squared norm:
\[
\vec w^\top(\mathbb{X}^\top\mathbb{X})\vec w \;=\; (\mathbb{X}\vec w)^\top(\mathbb{X}\vec w) \;=\; \lVert \mathbb{X}\vec w\rVert^2 \;\ge\; 0.
\]
Same trick for $\mathbb{X}\mathbb{X}^\top$: $\vec u^\top(\mathbb{X}\mathbb{X}^\top)\vec u = \lVert\mathbb{X}^\top\vec u\rVert^2 \ge 0$ for any $\vec u\in\mathbb{R}^N$. Combined with the spectral theorem (boxed in the problem: real symmetric $\Rightarrow$ real eigenvalues), every eigenvalue of $\mathbb{X}^\top\mathbb{X}$ and $\mathbb{X}\mathbb{X}^\top$ is a real number $\ge 0$, so $\sigma_i := \sqrt{\lambda_i}$ is well-defined.

\paragraph{Part (ii) --- shared nonzero spectra [3 points].} The proof is a single algebraic identity. Suppose $\mathbb{X}^\top\mathbb{X}\,\vec v = \lambda\vec v$ with $\lambda \ne 0$. Left-multiply by $\mathbb{X}$:
\[
\mathbb{X}\bigl(\mathbb{X}^\top\mathbb{X}\,\vec v\bigr) \;=\; \lambda\,(\mathbb{X}\vec v) \;\Longleftrightarrow\; \mathbb{X}\mathbb{X}^\top\,(\mathbb{X}\vec v) \;=\; \lambda\,(\mathbb{X}\vec v).
\]
This says $\mathbb{X}\vec v$ is an eigenvector of $\mathbb{X}\mathbb{X}^\top$ with the same eigenvalue $\lambda$, \emph{provided $\mathbb{X}\vec v \ne \vec 0$}.

\emph{Why} $\mathbb{X}\vec v\ne\vec 0$ \emph{here.} If $\mathbb{X}\vec v = \vec 0$ then $\mathbb{X}^\top\mathbb{X}\,\vec v = \mathbb{X}^\top\vec 0 = \vec 0 = \lambda\vec v$, which forces $\lambda = 0$ (since $\vec v$ is a nonzero eigenvector) and contradicts $\lambda\ne 0$. So as long as we restrict to \emph{nonzero} eigenvalues, $\mathbb{X}\vec v$ is automatically nonzero --- exactly why the statement is ``shared \emph{nonzero} eigenvalues.''

The reverse direction is identical: if $\mathbb{X}\mathbb{X}^\top\vec u = \lambda\vec u$ with $\lambda \ne 0$, then $\mathbb{X}^\top\vec u$ is an eigenvector of $\mathbb{X}^\top\mathbb{X}$ with the same $\lambda$. Defining $\sigma_i = \sqrt{\lambda_i}$ gives the clean statement: the nonzero eigenvalues of $\mathbb{X}^\top\mathbb{X}$ and $\mathbb{X}\mathbb{X}^\top$ coincide, both equal to $\sigma_i^2$. The unit right singular vectors $\vec v_i$ (eigenvectors of $\mathbb{X}^\top\mathbb{X}$) and the unit left singular vectors $\vec u_i := \mathbb{X}\vec v_i/\sigma_i$ (eigenvectors of $\mathbb{X}\mathbb{X}^\top$) are paired through exactly this construction.

\begin{insight}
\textbf{Geometric reading of the SVD.} Assembling 2d into one statement: every $N\times M$ matrix factors as $\mathbb{X} = \mathbb{U}\Sigma\mathbb{V}^\top$, where $\mathbb{U}$ and $\mathbb{V}$ are orthogonal (their columns are the left and right singular vectors) and $\Sigma$ is diagonal with the singular values. Geometrically, $\mathbb{X}$ acts on space as \emph{rotate} ($\mathbb{V}^\top$) $\to$ \emph{stretch by $\sigma_i$ along each new axis} ($\Sigma$) $\to$ \emph{rotate into output space} ($\mathbb{U}$). The unit sphere is sent to an ellipsoid whose semi-axes have lengths $\sigma_1\ge\sigma_2\ge\dots\ge\sigma_r > 0$.
\end{insight}

\paragraph{Forward reference to Q3.} The covariance matrix from Q3 is $\mathbb{C} = \mathbb{X}^\top\mathbb{X}/\tilde N$ (feature--feature, $M\times M$) and the Gram matrix is $\mathbb{G} = \mathbb{X}\mathbb{X}^\top/\tilde N$ (sample--sample, $N\times N$). The shared-nonzero-spectrum result just proved is \emph{why} PCA can be done in either space and recover the same principal components; the choice between them in Q3 is exactly the cost argument that 2e demands you justify.

\section*{2e) Power Method for SVD: a complete worked example [4 points]}

\paragraph{Goal.} (i) Write language-agnostic \textbf{pseudocode} that takes an arbitrary $N\times M$ matrix $\mathbb{X}$ and returns $\sigma_1, \vec v_1, \vec u_1$ using only the Power Method and matrix--vector products; justify the choice between $\mathbb{X}^\top\mathbb{X}$ and $\mathbb{X}\mathbb{X}^\top$ on cost grounds. (ii) Apply the algorithm \emph{by hand} (no computer) to a small $\mathbb{X}$ and verify $\mathbb{X}\vec v_1 = \sigma_1\vec u_1$ and $\mathbb{X}^\top\vec u_1 = \sigma_1\vec v_1$.

\paragraph{Part (i) --- pseudocode [2 points].} Two things earn the marks here: the algorithm itself, and the design choice of which symmetric matrix to feed it. Read the constraint first, then the cost argument, then the listing.

\begin{warn}
\textbf{Hard constraint.} The problem's box on this part says: \emph{``pseudocode only --- no Python, no other programming language.''} Use plain mathematical notation ($\mathbb{X}^\top\vec y$, $\lVert\vec y\rVert$, ``while $\dots$ repeat''), \textbf{not} Python syntax such as \texttt{@}, \texttt{np.\allowbreak linalg.\allowbreak norm}, or \texttt{for k \allowbreak in \allowbreak range(K)}.
\end{warn}

\emph{Design choice (where the cost-argument mark lives).} The Power Method needs a symmetric input.
\begin{itemize}[leftmargin=*]
\item Apply it to $\mathbb{X}^\top\mathbb{X}$ ($M\times M$): forming the matrix once costs $O(NM^2)$ flops and $M^2$ storage.
\item Apply it to $\mathbb{X}\mathbb{X}^\top$ ($N\times N$): same analysis with $N\leftrightarrow M$.
\end{itemize}
Pick the \textbf{smaller} side: $\mathbb{X}^\top\mathbb{X}$ when $M \le N$, $\mathbb{X}\mathbb{X}^\top$ when $N < M$. For Q3 ($N=1000$, $M=3041$) the Gram matrix is $1000\times 1000$ versus the $3041\times 3041$ covariance --- nine times cheaper to store.

\emph{Further trick: don't materialise the product.} You never need to form $\mathbb{X}^\top\mathbb{X}$ explicitly: each Power-Method step needs only $(\mathbb{X}^\top\mathbb{X})\vec v$, which you can compute as $\mathbb{X}^\top(\mathbb{X}\vec v)$ in two matrix--vector products of cost $O(NM)$ each, with no $M\times M$ storage. For very thin tall matrices the saving is enormous and is worth a one-line mention in the report.

\medskip
\noindent\textbf{Algorithm (pseudocode).} Inputs: matrix $\mathbb{X}\in\mathbb{R}^{N\times M}$; tolerance $\mathrm{tol}>0$; iteration cap $K$. Outputs: $\sigma_1, \vec v_1, \vec u_1$.

\begin{enumerate}[label=\textbf{\arabic*.},leftmargin=*]
\item \textbf{Choose the smaller side.} If $M \le N$, work in $\mathbb{R}^M$ and apply $\mathbb{A} := \mathbb{X}^\top\mathbb{X}$ implicitly via $\vec y \mapsto \mathbb{X}^\top(\mathbb{X}\vec y)$. Otherwise work in $\mathbb{R}^N$ and apply $\mathbb{A} := \mathbb{X}\mathbb{X}^\top$ implicitly via $\vec y \mapsto \mathbb{X}(\mathbb{X}^\top\vec y)$.
\item \textbf{Initialise.} Draw a random nonzero $\vec v_0$ in the chosen space; normalise: $\vec v_0 \leftarrow \vec v_0 / \lVert\vec v_0\rVert$.
\item \textbf{Power iteration.} For $k = 0, 1, 2, \dots, K-1$:
\begin{enumerate}[label=(\alph*),leftmargin=*]
\item Compute $\vec w := \mathbb{A}\vec v_k$ via two matrix--vector products (never the explicit product matrix).
\item Normalise: $\vec v_{k+1} := \vec w / \lVert\vec w\rVert$.
\item If $1 - |\vec v_{k+1}\cdot\vec v_k| < \mathrm{tol}$, break (angle convergence; insensitive to a $\pm$ sign flip between steps).
\end{enumerate}
\item \textbf{Read off the singular value} via the Rayleigh quotient: $\sigma_1^2 := \vec v_{k+1}^{\,\top}\,\mathbb{A}\vec v_{k+1}$, then $\sigma_1 := \sqrt{\sigma_1^2}$.
\item \textbf{Companion vector via one matrix--vector product.} If iterated in $\mathbb{R}^M$, set $\vec v_1 := \vec v_{k+1}$ and $\vec u_1 := \mathbb{X}\vec v_1/\sigma_1$. If iterated in $\mathbb{R}^N$, set $\vec u_1 := \vec v_{k+1}$ and $\vec v_1 := \mathbb{X}^\top\vec u_1/\sigma_1$.
\item \textbf{Return} $(\sigma_1, \vec v_1, \vec u_1)$.
\end{enumerate}

\paragraph{Part (ii) --- hand calculation [2 points].} Apply the algorithm to
\[
\mathbb{X} \;=\; \begin{pmatrix} 1 & 0 \\ 0 & 1 \\ 1 & 1 \end{pmatrix} \quad (N=3,\ M=2).
\]
Since $M=2<N=3$, iterate on $\mathbb{X}^\top\mathbb{X}\in\mathbb{R}^{2\times 2}$. Form it explicitly (small enough by hand):
\[
\mathbb{X}^\top \;=\; \begin{pmatrix} 1 & 0 & 1 \\ 0 & 1 & 1 \end{pmatrix},
\qquad
\mathbb{X}^\top\mathbb{X} \;=\; \begin{pmatrix} 2 & 1 \\ 1 & 2 \end{pmatrix}.
\]
\emph{Eigenvalues.} $\det(\mathbb{X}^\top\mathbb{X} - \lambda I) = (2-\lambda)^2 - 1 = 0 \Rightarrow \lambda \in \{3,\,1\}$. Dominant $\lambda_1 = 3$, so $\sigma_1 = \sqrt{3}$.

\emph{Dominant eigenvector} $\vec v_1$. From $(\mathbb{X}^\top\mathbb{X} - 3I)\vec v = \vec 0$,
\[
\begin{pmatrix} -1 & 1 \\ 1 & -1 \end{pmatrix}\begin{pmatrix} a \\ b \end{pmatrix} = \vec 0 \;\Longrightarrow\; b = a, \quad \vec v_1 \;=\; \frac{1}{\sqrt{2}}\begin{pmatrix} 1 \\ 1 \end{pmatrix}.
\]
(A genuine Power-Method run would iterate $\vec v_k \mapsto \mathbb{X}^\top(\mathbb{X}\vec v_k)/\lVert\cdot\rVert$ from a random unit start and lock onto this direction; $M=2$ is small enough to solve the characteristic polynomial directly, which the problem explicitly permits for the hand step.)

\emph{Companion left singular vector} $\vec u_1$.
\[
\vec u_1 \;=\; \frac{\mathbb{X}\vec v_1}{\sigma_1} \;=\; \frac{1}{\sqrt{3}}\cdot\frac{1}{\sqrt{2}}\begin{pmatrix} 1 \\ 1 \\ 2 \end{pmatrix} \;=\; \frac{1}{\sqrt{6}}\begin{pmatrix} 1 \\ 1 \\ 2 \end{pmatrix}.
\]

\begin{sanity}
\textbf{Verification of} $\mathbb{X}\vec v_1 = \sigma_1\vec u_1$.
\[
\mathbb{X}\vec v_1 = \frac{1}{\sqrt{2}}\begin{pmatrix} 1 \\ 1 \\ 2 \end{pmatrix},
\qquad
\sigma_1\vec u_1 = \sqrt{3}\cdot \frac{1}{\sqrt{6}}\begin{pmatrix} 1 \\ 1 \\ 2 \end{pmatrix} = \frac{1}{\sqrt{2}}\begin{pmatrix} 1 \\ 1 \\ 2 \end{pmatrix}. \quad\checkmark
\]
\textbf{Verification of} $\mathbb{X}^\top\vec u_1 = \sigma_1\vec v_1$.
\[
\mathbb{X}^\top\vec u_1 = \frac{1}{\sqrt{6}}\begin{pmatrix} 1+0+2 \\ 0+1+2 \end{pmatrix} = \frac{3}{\sqrt{6}}\begin{pmatrix} 1 \\ 1 \end{pmatrix},
\qquad
\sigma_1\vec v_1 = \frac{\sqrt{3}}{\sqrt{2}}\begin{pmatrix} 1 \\ 1 \end{pmatrix} = \frac{3}{\sqrt{6}}\begin{pmatrix} 1 \\ 1 \end{pmatrix}. \quad\checkmark
\]
$\mathbb{X}$ stretches $\vec v_1$ by $\sigma_1=\sqrt{3}$ into the direction $\vec u_1$, and $\mathbb{X}^\top$ does the reverse.
\end{sanity}

\begin{warn}
\textbf{Pitfall.} It is tempting to obtain $\vec u_1$ by eigendecomposing $\mathbb{X}\mathbb{X}^\top$ independently. Both approaches yield an eigenvector of $\mathbb{X}\mathbb{X}^\top$ with eigenvalue $\sigma_1^2$, but the second can return $-\vec u_1$ instead of $+\vec u_1$, in which case the verification $\mathbb{X}\vec v_1 = \sigma_1\vec u_1$ will appear to fail by a sign. Use $\vec u_1 := \mathbb{X}\vec v_1/\sigma_1$ to pin the sign automatically.
\end{warn}

\paragraph{What goes in the report.} For the 10-page submission, 2e should be terse: the pseudocode listing (boxed), a one-line cost justification, then the hand calculation as five labelled equations (form $\mathbb{X}^\top\mathbb{X} \to$ eigenvalues $\to \vec v_1 \to \vec u_1 \to$ both verifications). Do not narrate Power-Method iterations themselves; the marks here are for the algorithm, the cost justification, and the verifications, not for showing several iterations converge to the same answer.

\section*{Summary: how Q1 leads into Q2 and Q3}

\begin{itemize}[leftmargin=*]
    \item \textbf{Q1's lesson} is that the long-time behaviour of a linear system is read off from the eigenvector with the largest eigenvalue. The other eigenvectors describe transient deviations that geometrically decay.
    \item \textbf{Q2 generalises this} into the \emph{Power Method}: when the dominant eigenvalue is unique, iterating any vector by the matrix converges to that eigenvector. Q2 also asks you to make the convergence rate quantitative (it is the spectral gap $|\lambda_2/\lambda_1|$ from Q1), and to extend the idea to rectangular matrices via SVD.
    \item \textbf{Q3 applies the same machinery to a real $1000\times 3041$ infrared spectroscopy dataset}: the dominant eigenvectors of the covariance matrix are the directions of greatest spectroscopic variation, which separate the unknown samples into clusters. Q1's Markov example is the toy that makes this principle visible; Q3 is where it pays off.
\end{itemize}

\vfill
\hrule
\smallskip
\noindent\textit{End of guide. The actual report should be much terser than this document --- aim for a few clean lines per part. This guide is meant to make the underlying ideas obvious, not to be a template for length.}

\end{document}
